\paragraph{端点层析原理：由单径向模长恢复“轴外零点”全数据的可审计反演链}\label{par:xi-endpoint-tomography-principle}
在 Hardy/Smirnov 条件口径下，圆盘化完成化函数 \(F_{\zeta}\) 具有内外分解 \(F_{\zeta}=BSO\)（附录 \S\ref{subsec:unit-disk-inner-outer}）。
其中 Blaschke 因子 \(B\) 的盘内零点集合 \(\{a_k\}\subset\mathbb{D}\) 编码“离临界线偏移”的全部离散自由度；其端点 \(-1\) 的吸收系数
\(\Phi_{-1}(B)=\sum_k P_{a_k}(-1)\)
由单径向一阶衰减率读出（命题 \ref{prop:xi-endpoint-absorption-coefficient}，并见推论 \ref{cor:xi-endpoint-absorption-additive-tomography}）。
在无奇异内因子 \(S\equiv 1\) 且有限 Blaschke 的情形，该量亦等于端点相位电流密度 \(j_{\mathrm{in}}(\pi)\)（见 \eqref{eq:app-endpoint-phase-current-lowerbound}）。
本段给出一个更强的层析化结果：在有限缺陷（有限 Blaschke）口径下，\emph{整条端点径向模长剖面}
\(
r\mapsto |B(-r)|
\)
不仅给出 \(\sum_k\delta_k\) 的总量，而且携带 \((\gamma_k,\delta_k)\) 的全数据，并在解析延拓意义下具有反演唯一性。

\begin{definition}[端点径向吸收剖面]\label{def:xi-endpoint-radial-absorption-profile}
令 \(B\) 为有限 Blaschke 乘积。
定义端点 \(-1\) 的\textbf{径向吸收剖面}为
\begin{equation}\label{eq:xi-endpoint-radial-absorption-profile}
\boxed{\;
\mathcal{A}_B(r):=-\frac{\dd}{\dd r}\log|B(-r)|
\qquad(0<r<1).\;}
\end{equation}
\end{definition}

\begin{theorem}[端点径向吸收的全核表示：一维剖面携带二维零点数据]\label{thm:xi-endpoint-absorption-kernel-representation}
令 \(B\) 为有限 Blaschke 乘积，零点多重集为 \(\{a_k\}_{k=1}^{N}\subset\mathbb{D}\)（按重数计）。
令
\[
t_k:=\mathcal{C}^{-1}(a_k)=\gamma_k+\mathrm{i}\delta_k\in\mathbb{H},
\qquad (\gamma_k\in\RR,\ \delta_k>0),
\]
其中 \(\mathcal{C}\) 为本文固定的 Cayley 紧化（从 \(\mathbb{H}\) 到 \(\mathbb{D}\)，见命题 \ref{prop:app-busemann-poisson-minus1} 与推论 \ref{cor:app-offcritical-horocycle-busemann}）。
则对任意 \(r\in(0,1)\) 有显式核分解
\begin{equation}\label{eq:xi-endpoint-absorption-kernel-sum}
\boxed{\;
\mathcal{A}_B(r)
=\sum_{k=1}^{N}
\frac{4\delta_k}{(1-r)^2\gamma_k^2+\bigl(1+r+(1-r)\delta_k\bigr)^2}\; }.
\end{equation}
并且端点极限满足
\begin{equation}\label{eq:xi-endpoint-absorption-kernel-limit}
\boxed{\;
\lim_{r\uparrow 1}\mathcal{A}_B(r)
=\sum_{k=1}^{N}\delta_k
=\sum_{k=1}^{N}P_{a_k}(-1)
=\Phi_{-1}(B)\; }.
\end{equation}
\end{theorem}

\begin{proof}
将 \(B\) 写作 Blaschke 因子乘积 \(B=\prod_{k=1}^{N}b_{a_k}\)（忽略常相位因子）。
对单因子 \(b_a(w)=\frac{a-w}{1-\overline{a}\,w}\) 在实径向 \(w=-r\) 上有
\[
\frac{\dd}{\dd r}\log|b_a(-r)|
=\Re\!\left(\frac{1}{a+r}-\frac{\overline a}{1+r\overline a}\right),
\]
从而 \(\mathcal{A}_B(r)=-\sum_k \frac{\dd}{\dd r}\log|b_{a_k}(-r)|\)。
\par
令 \(a=\mathcal{C}(t)\) 且 \(t=\gamma+\mathrm{i}\delta\in\mathbb{H}\)。
由 \(\mathcal{C}\) 的闭式表达（等价于命题 \ref{prop:app-busemann-poisson-minus1} 中 \(B^{-1}(w)=\Im(\mathcal{C}^{-1}(w))\) 的 Cayley 选取）可将上式分母通分并整理为
\[
-\frac{\dd}{\dd r}\log|b_a(-r)|
=\frac{4\delta}{(1-r)^2\gamma^2+\bigl(1+r+(1-r)\delta\bigr)^2}.
\]
对 \(k=1,\dots,N\) 求和即得 \eqref{eq:xi-endpoint-absorption-kernel-sum}。
\par
令 \(r\uparrow 1\) 则分母趋于 \(4\)，故每一项趋于 \(\delta_k\)，从而
\(\lim_{r\uparrow 1}\mathcal{A}_B(r)=\sum_k\delta_k\)。
再由 \(P_{a_k}(-1)=B^{-1}(a_k)=\Im(\mathcal{C}^{-1}(a_k))=\delta_k\)（\eqref{eq:app-poisson-kernel-endpoint-minus1}）与 \(\Phi_{-1}(B)=\sum_k P_{a_k}(-1)\) 的定义（命题 \ref{prop:xi-endpoint-absorption-coefficient}）得到 \eqref{eq:xi-endpoint-absorption-kernel-limit}。证毕。
\end{proof}

\leanverified{paper\_xi\_endpoint\_absorption\_kernel\_representation}

\begin{theorem}[端点喷射不变量：\texorpdfstring{$\mathcal{A}_B$}{A\_B} 在 \texorpdfstring{$r=1$}{r=1} 的低阶导数给出可核验“矩”]\label{thm:xi-endpoint-absorption-jet-moments}
在定理 \ref{thm:xi-endpoint-absorption-kernel-representation} 的有限 Blaschke 设定下，函数 \(\mathcal{A}_B\) 在 \(r\uparrow 1\) 的邻域光滑，并满足端点喷射恒等式
\[
\boxed{\;
\mathcal{A}_B(1)=\sum_{k}\delta_k,\qquad
\mathcal{A}_B'(1)=\sum_{k}\delta_k(\delta_k-1),\qquad
\mathcal{A}_B''(1)=\frac12\sum_{k}\delta_k\bigl(3(\delta_k-1)^2-\gamma_k^2\bigr).\;}
\]
当 \(\delta_k\in(0,\tfrac12)\)（离临界线偏移口径）时，有 \(\mathcal{A}_B'(1)<0\)；并且二阶项中
\(-\tfrac12\sum_k\delta_k\gamma_k^2\)
给出对大高度 \(|\gamma|\) 的二阶抑制响应。
\end{theorem}

\begin{proof}
由 \eqref{eq:xi-endpoint-absorption-kernel-sum}，\(\mathcal{A}_B\) 为有限项的有理函数和，故在 \(r=1\) 的一侧邻域光滑。
对单项核
\(
K_{(\gamma,\delta)}(r)=\frac{4\delta}{(1-r)^2\gamma^2+(1+r+(1-r)\delta)^2}
\)
作 \(u=1-r\) 的展开即可得
\[
K_{(\gamma,\delta)}(1)=\delta,\quad
K_{(\gamma,\delta)}'(1)=\delta(\delta-1),\quad
K_{(\gamma,\delta)}''(1)=\tfrac12\,\delta\bigl(3(\delta-1)^2-\gamma^2\bigr),
\]
对 \(k\) 求和即得结论。证毕。
\end{proof}

\begin{theorem}[有限缺陷情形的端点反演唯一性：\texorpdfstring{$\mathcal{A}_B$}{A\_B} 的单变量数据决定零点多重集]\label{thm:xi-endpoint-absorption-inversion-uniqueness}
在定理 \ref{thm:xi-endpoint-absorption-kernel-representation} 的设定下，令
\[
q:=\frac{1-r}{1+r}\in(0,1),
\qquad r=\frac{1-q}{1+q}.
\]
则 \(q\mapsto \mathcal{A}_B\!\bigl(\frac{1-q}{1+q}\bigr)\) 在 \(q\) 的复平面上为有理函数，其复极点集合由 \(\{t_k=\gamma_k+\mathrm{i}\delta_k\}\) 完全决定：对每个 \(k\) 存在一对共轭极点
\begin{equation}\label{eq:xi-endpoint-absorption-q-poles}
\boxed{\;
q_{k,\pm}=\frac{-\delta_k\pm \mathrm{i}\gamma_k}{\delta_k^2+\gamma_k^2}
=-\frac{\delta_k\mp \mathrm{i}\gamma_k}{\delta_k^2+\gamma_k^2}\;}
\end{equation}
且相应留数仅依赖于 \(\delta_k\)。
因此：
\begin{enumerate}
  \item 若两有限 Blaschke 乘积 \(B_1,B_2\) 满足 \(\mathcal{A}_{B_1}(r)=\mathcal{A}_{B_2}(r)\) 在某个非平凡区间上成立，则其零点多重集在 Cayley 坐标下必一致（按重数计）：
  \[
  \{\mathcal{C}^{-1}(a_k^{(1)})\}=\{\mathcal{C}^{-1}(a_k^{(2)})\}.
  \]
  \item 特别地，端点径向模长剖面 \(r\mapsto |B(-r)|\)（等价于 \(\mathcal{A}_B\)）在有限缺陷口径下决定全部参数 \((\gamma_k,\delta_k)\)，至多相差于由共轭对称与排列导致的自然不可辨识；该不可辨识与零点共轭配对/自对偶对称完全相容。
\end{enumerate}
\end{theorem}

\begin{proof}
将 \eqref{eq:xi-endpoint-absorption-kernel-sum} 代入 \(r=\frac{1-q}{1+q}\) 并用恒等式
\(
1-r=\frac{2q}{1+q},\ 1+r=\frac{2}{1+q}
\)
整理得
\[
\mathcal{A}_B\!\Bigl(\frac{1-q}{1+q}\Bigr)
=\sum_{k=1}^{N}
\frac{\delta_k(1+q)^2}{q^2\gamma_k^2+(1+q\delta_k)^2},
\]
从而为有理函数。其极点由分母的零点给出：
\(
q^2\gamma_k^2+(1+q\delta_k)^2=0
\)
等价于
\(
(\delta_k^2+\gamma_k^2)q^2+2\delta_k q+1=0
\)，
解得 \eqref{eq:xi-endpoint-absorption-q-poles}，并由共轭对称得到成对极点。
留数只依赖于 \(\delta_k\) 为直接代入计算。
\par
若 \(\mathcal{A}_{B_1}=\mathcal{A}_{B_2}\) 在一段区间上相等，则二者在 \(q\)-变量下对应的有理函数在开集上相等，从而由解析延拓（恒等定理）知其作为亚纯函数处处相等，极点集合与重数必须一致；于是 \(\{(\gamma_k,\delta_k)\}\)（等价地 \(\{t_k\}\)）按重数一致，结论 (1) 成立。
结论 (2) 由 \(\mathcal{A}_B=-\frac{\dd}{\dd r}\log|B(-r)|\) 与积分恢复 \(|B(-r)|\)（结合 \(B(0)\neq 0\) 的归一化）推出。证毕。
\end{proof}

\begin{corollary}[端点剖面系数的 C-finite 刚性：常系数递推与 Hankel 秩证书]\label{cor:xi-endpoint-profile-cfinite-hankel-rank}
\leanverified{paper\_xi\_endpoint\_profile\_cfinite\_hankel\_rank}
在定理 \ref{thm:xi-endpoint-absorption-inversion-uniqueness} 的设定下，令
\[
F(q):=\mathcal{A}_B\!\Bigl(\frac{1-q}{1+q}\Bigr)=\sum_{n\ge 0}c_n q^n\qquad(|q|\ \text{充分小}).
\]
则 \(F\) 为有理函数，并可取分母
\[
\boxed{\ D(q):=\prod_{k=1}^{N}\bigl((\delta_k^2+\gamma_k^2)q^2+2\delta_k q+1\bigr)\in\ZZ[q]\ }.
\]
因此系数序列 \(\{c_n\}_{n\ge 0}\) 为一个严格的 C-finite 序列：存在阶 \(2N\) 的常系数线性递推使对所有充分大 \(n\) 有
\[
c_{n+2N}=\sum_{j=0}^{2N-1}\alpha_j\,c_{n+j}\qquad(\alpha_j\in\QQ).
\]
若进一步假设极点对 \(\{q_{k,\pm}\}\) 两两不同（等价于 \(\{t_k=\gamma_k+\mathrm{i}\delta_k\}\) 在 \(\mathbb{H}\) 中按重数互异），则上述阶数为最小阶；等价地，Hankel 矩阵
\(
\mathsf{H}=(c_{i+j})_{0\le i,j\le 2N}
\)
的秩恰为 \(2N\)，从而 \(N\) 可由有限截断的线性代数证书精确读出。
\end{corollary}
\begin{proof}
由定理 \ref{thm:xi-endpoint-absorption-inversion-uniqueness} 的有理表达式，\(F\) 的极点恰为二次因子
\(
(\delta_k^2+\gamma_k^2)q^2+2\delta_k q+1
\)
的根 \(q_{k,\pm}\)。取 \(D(q)\) 为这些二次因子的乘积即可使 \(D(q)F(q)\) 成为多项式，从而 \(F\) 为有理函数且其 Taylor 系数满足由 \(D\) 的系数给出的常系数线性递推（将 \(D(q)F(q)\) 的幂级数系数逐项比较即可）。
\par
当极点两两不同且无因子消去时，\(D\) 为最小分母（up to scalar），其次数 \(2N\) 即最小递推阶；而最小递推阶与 Hankel 秩刻画等价。证毕。
\end{proof}

\begin{corollary}[端点数据的 Prony 可逆性：\texorpdfstring{$4N$}{4N} 个系数恢复全部 \texorpdfstring{$(\gamma_k,\delta_k)$}{(gamma,delta)}]\label{cor:xi-endpoint-profile-prony-invertibility}
\leanverified{paper\_xi\_endpoint\_profile\_prony\_invertibility}
在定理 \ref{thm:xi-endpoint-absorption-inversion-uniqueness} 的设定下，若进一步假设极点对 \(\{q_{k,\pm}\}\) 两两不同，则给定 \(F(q)\) 在 \(q=0\) 的任意 \(4N\) 个连续 Taylor 系数 \(c_0,c_1,\dots,c_{4N-1}\)，即可在原则上唯一恢复参数多重集 \(\{(\gamma_k,\delta_k)\}_{k=1}^{N}\)（至多相差于共轭配对与排列的自然不可辨识）。
\end{corollary}
\begin{proof}
由推论 \ref{cor:xi-endpoint-profile-cfinite-hankel-rank}，\(F\) 的最小分母 \(D(q)\) 次数为 \(2N\)，从而 \(\{c_n\}\) 满足最小阶 \(2N\) 递推；该递推（等价的 annihilating 多项式）可由 Hankel 线性系统/Prony 核空间法（亦可用 Berlekamp--Massey）从 \(c_0,\dots,c_{4N-1}\) 唯一恢复，进而得到 \(D(q)\)（归一化首项为 \(1\)）。
将 \(D\) 分解为二次因子即可恢复全部极点 \(q_{k,\pm}\)；再由 \eqref{eq:xi-endpoint-absorption-q-poles} 的和积关系可解出 \((\gamma_k,\delta_k)\)。
最后，留数（等价于各项系数）可由将 \(F\) 写作部分分式并匹配前 \(2N\) 个系数的线性系统唯一确定。证毕。
\end{proof}

\begin{corollary}[对 \texorpdfstring{$F_\zeta$}{F\_zeta} 的端点证书链：从单径向剖面到离线零点存在性]\label{cor:xi-endpoint-tomography-certificate-chain}
在附录 \S\ref{subsec:unit-disk-inner-outer} 的因子化口径下，若进一步假设 \(S\equiv 1\)（无奇异内因子），则端点径向剖面 \(r\mapsto |B(-r)|\) 由 \(F_\zeta\) 的端点径向模长剖面与外因子剖面共同决定：
\[
\log|F_\zeta(-r)|=\log|B(-r)|+\log|O(-r)|.
\]
因此，在外因子 \(O\) 的径向对数导数可被可审计地控制/剥离的条件下，\(\mathcal{A}_B\) 的非零性给出离线零点存在性的端点证书；并且在有限缺陷近似（有限 Blaschke）下，\(\mathcal{A}_B\) 甚至可反演出全部 \((\gamma_k,\delta_k)\)，从而把“离线零点”从二维隐数据转写为单变量剖面函数的可审计谱指纹（定理 \ref{thm:xi-endpoint-absorption-inversion-uniqueness}）。
\end{corollary}

\begin{theorem}[无穷 Blaschke 漂移剖面的亚纯结构、可识别性与局部稳定]\label{thm:xi-endpoint-infinite-blaschke-drift-identifiability-stability}
令 \(B\) 为（可能无穷）Blaschke 乘积，并记 Cayley 参数
\(
t_k=\gamma_k+\mathrm{i}\delta_k\in\mathbb H
\)（\(\delta_k>0\)）。
假设可和条件
\[
\sum_k\frac{\delta_k}{\gamma_k^2+\delta_k^2}<\infty.
\]
定义漂移剖面
\[
\mathcal A_B(r):=-\frac{\dd}{\dd r}\log|B(-r)|,\qquad r\in(0,1),
\]
并令
\[
F_B(q):=\mathcal A_B\!\left(\frac{1-q}{1+q}\right).
\]
则成立：
\begin{enumerate}
  \item \(F_B\) 在 \(\Re q>0\) 解析，并可唯一亚纯延拓到
  \(
  \CC\setminus\{q_{k,\pm}\}
  \)，其中
  \[
  q_{k,\pm}=\frac{-\delta_k\pm \mathrm{i}\gamma_k}{\delta_k^2+\gamma_k^2}.
  \]
  \item 若 \(B_1,B_2\) 满足同样可和条件，且在任一具有聚点的集合
  \(I\subset(0,1)\) 上有
  \(
  \mathcal A_{B_1}(r)=\mathcal A_{B_2}(r)
  \)，
  则二者极点多重集一致，等价于
  \(
  \{(\gamma_k,\delta_k)\}
  \)
  一致（仅差排列与共轭配对）。
  \item 设 \(F_B\) 在闭圆盘 \(\overline{D(q_0,\eta)}\) 内仅有一个简单极点 \(q_0\)，留数 \(\alpha\neq 0\)；
  设 \(\widetilde F\) 在该盘内亦仅有一个简单极点 \(\widetilde q_0\)，并满足
  \[
  \sup_{|q-q_0|=\eta}|F_B(q)-\widetilde F(q)|\le \varepsilon
  \quad\text{且}\quad
  \varepsilon<\frac{|\alpha|}{2\eta}.
  \]
  则有显式稳定界
  \[
  |\widetilde q_0-q_0|
  \le
  \frac{2\eta^2}{|\alpha|}\,\varepsilon.
  \]
\end{enumerate}
\end{theorem}

\begin{proof}
由定理 \ref{thm:xi-endpoint-absorption-kernel-representation} 的单因子公式，
\[
F_B(q)=\sum_k \frac{\delta_k(1+q)^2}{q^2\gamma_k^2+(1+q\delta_k)^2}.
\]
对任意紧集 \(K\Subset\{q:\Re q>0\}\)，每项可由常数 \(C_K\) 估计为
\[
\left|\frac{\delta_k(1+q)^2}{q^2\gamma_k^2+(1+q\delta_k)^2}\right|
\le
C_K\frac{\delta_k}{\gamma_k^2+\delta_k^2},
\]
而右端可和，故由 Weierstrass 判别法得在 \(\Re q>0\) 上局部一致收敛并解析。
同理在任意避开 \(\{q_{k,\pm}\}\) 的紧集上局部一致收敛，得到唯一亚纯延拓。

若 \(B_1,B_2\) 在 \(I\subset(0,1)\) 上剖面相等，则对应 \(F_{B_1},F_{B_2}\) 在 \(\Re q>0\) 的带聚点集合上相等，由恒等定理得该半平面内恒等；再由亚纯延拓唯一性，极点多重集一致。
由
\[
q_{k,+}=\frac{-\delta_k+\mathrm{i}\gamma_k}{\delta_k^2+\gamma_k^2}
\quad\Longrightarrow\quad
\delta_k=\frac{-\Re q_{k,+}}{|q_{k,+}|^2},\ 
\gamma_k=\frac{\Im q_{k,+}}{|q_{k,+}|^2},
\]
可恢复参数多重集（到排列/共轭配对）。

最后证明局部稳定界。记
\[
\alpha=\frac{1}{2\pi i}\oint_{|q-q_0|=\eta}F_B(q)\,\dd q,\qquad
\widetilde\alpha=\frac{1}{2\pi i}\oint_{|q-q_0|=\eta}\widetilde F(q)\,\dd q .
\]
由边界误差界，
\[
|\widetilde\alpha-\alpha|
\le
\frac{1}{2\pi}\cdot(2\pi\eta)\varepsilon
=\eta\varepsilon
<\frac{|\alpha|}{2},
\]
从而 \(|\widetilde\alpha|\ge |\alpha|/2\)。
又由于 \((q-q_0)F_B\) 在盘内全纯（唯一极点已消去），有
\[
\frac{1}{2\pi i}\oint_{|q-q_0|=\eta}(q-q_0)\widetilde F(q)\,\dd q
=
\widetilde\alpha(\widetilde q_0-q_0)
=
\frac{1}{2\pi i}\oint_{|q-q_0|=\eta}(q-q_0)\bigl(\widetilde F-F_B\bigr)\,\dd q.
\]
故
\[
|\widetilde\alpha|\,|\widetilde q_0-q_0|
\le
\frac{1}{2\pi}\cdot(2\pi\eta)\cdot\eta\varepsilon
=\eta^2\varepsilon.
\]
代入 \(|\widetilde\alpha|\ge |\alpha|/2\) 即得
\[
|\widetilde q_0-q_0|
\le
\frac{2\eta^2}{|\alpha|}\,\varepsilon.
\]
证毕。
\end{proof}

\endinput
